\section{Review of Existing Techniques}
\label{sec:related}
CLIP learns a shared image--text representation, allowing class descriptions to act as a classifier without downstream training labels~\citep{radford2021clip}. Adaptation can change the text input, the representation, or the classifier. CoOp learns continuous context vectors while keeping the pretrained model fixed~\citep{zhou2022coop}. CLIP-Adapter instead learns a small feature transformation and blends it with the original representation~\citep{gao2024adapter}. Tip-Adapter retrieves information from a cache of labelled support features and combines it with zero-shot logits~\citep{zhang2022tip}. We evaluate linear probes and the training-free form of Tip-Adapter; prompt learning and a trainable feature adapter are related approaches, not additional methods evaluated here.

Preserving pretrained knowledge is also central to robust fine-tuning. Kumar et al.~\citep{kumar2022lpft} show that updating the encoder can distort useful features under distribution shift, motivating linear probing followed by fine-tuning. WiSE-FT interpolates zero-shot and fine-tuned model weights to improve the accuracy--robustness trade-off~\citep{wortsman2022wiseft}. Our encoder never changes, so its negative flips cannot be attributed to encoder feature distortion. Moreover, our fixed-mixing control combines output probabilities, rather than model weights. These distinctions matter when interpreting similarities to robust fine-tuning results.

Few-shot comparisons also depend on model selection. Silva-Rodríguez et al.~\citep{silva2024clap} find that strong results from existing adapters can rely on task-specific hyperparameters and substantial validation data. They revisit zero-shot-initialized linear probing and introduce class-adaptive constraints in CLAP. Their probe also normalizes weights after updates and uses augmentation; our A2 is a simpler, unconstrained implementation without either component. DomainBed emphasizes that the selection rule is part of the learning procedure~\citep{gulrajani2021domainbed}. We therefore select settings on other domains and explicitly count target support labels. This is few-shot domain adaptation, not the label-free target setting of domain generalization.

Sample-wise combination has several precedents. Maximum softmax probability provides a simple signal of whether a prediction is correct~\citep{hendrycks2017msp}. VRF weights zero-shot and fine-tuned predictions using distance to training examples corrected by fine-tuning~\citep{zhu2024vrf}; DaWin uses predictive entropy to choose sample-specific weight interpolation coefficients~\citep{oh2025dawin}. Our hard confidence gate and frozen-feature VRF variant test these ideas with very small support sets. They do not reproduce full VRF or DaWin. Fixed probability mixing is a necessary control because confidence affects the relative influence of its two probability vectors even when the mixing coefficient is constant.

Finally, negative-flip analysis predates this study. Positive-congruent training explicitly studies updates that turn correct predictions into errors~\citep{yan2021flips}. We use that definition to examine few-shot CLIP adaptation across support sizes, domains and methods, then compare attempted safeguards against matched random controls and fixed mixing. The contribution is this controlled measurement and diagnosis under a stated selection protocol, rather than a new definition of harm or a claim to have invented sample-wise combination.
